% Nouveau EGA V (cote 25): archive cote 25, site page 164 (PDF page 165).
% typescript. Draft compared: none.
% AI-transcribed by Claude Opus 5.5 (claude-opus-5-5, adaptive thinking, effort high), Message Batches API, batch ega25b (msgbatch_01A7n5YuXubUPNK1zsepi8RG), 2026-11-13; not yet checked by eye.
\scanpage{25/164}
\mspage{7}

k-isomorphe à $k[T]/\unsure{(}T^n)$ . Si $d=1$ ,\var{\inserted[typed]{x rat. sur k,}}{x rat. sur k,}{insertion}{ega25_25_164-v0001} la notion de point singulier
\var{\xout{ordinaire}}{}{typewriter strike-out}{ega25_25_164-v0002} élémentaire de multiplicité n correspond à la terminologie
classique \var{\xout{\illegible{1}}}{}{typewriter strike-out}{ega25_25_164-v0003}``point singulier ordinaire à n tangente\unsure{\var{\xout{s}}{}{typewriter strike-out}{ega25_25_164-v0004}|s} disctinctes'' .
\var{\margin{(}}{\placedmargin{(}}{marginal note placed}{ega25_25_164-v0005}(On aura \var{\inserted{déjà}}{déjà}{insertion}{ega25_25_164-v0006} explicité dans 15.2. que si x est un point de X rat sur k, alors
la notion de singularité élémentaire de mult. n, au sens absolu ou au
sens relatif, \var{\xout{\unsure{elle}}}{}{typewriter strike-out}{ega25_25_164-v0007} est la même; cela reste d'ailleurs valable si k(x)
est une extension finie séparable de k - sans doute d'ailleurs ``finie''
est inutile - et ce fait mérite d'être inséré en corollaire \var{\xout{\unsure{si on arrive}}}{}{typewriter strike-out}{ega25_25_164-v0008}.
ou en proposition\hand{)}.

NB J'avais inclus au par.20 un nº sur l'éclatement d'un
préschéma X le long d'un sous-préschéma \var{\inserted{fermé}}{fermé}{insertion}{ega25_25_164-v0009} Y tel que l'augmentation
$\underline{O}_X \longrightarrow \underline{O}_Y$ soit élémentaire de mult.n (auquel j'ai fait allusion dans
15.6.). Je signale qu'il s'agit d'un nº court et bien venu, dont les
seuls ingrédients sont les résultats généraux de lissité du par.17,et
la définition \var{\inserted{du début de}}{du début de}{insertion}{ega25_25_164-v0010} 17.15.\unsure{2}. \var{\inserted[typed]{(et encore pourrait-on éviter cette dernière).}}{(et encore pourrait-on éviter cette dernière).}{insertion}{ega25_25_164-v0011} Rien n'empêcherait d'incopporer ces résultats
dès ici, \var{\xout{\unsure{par dans}}}{}{typewriter strike-out}{ega25_25_164-v0012} par exemple immédiatement après la définition
\var{\xout{17}}{}{typewriter strike-out}{ega25_25_164-v0013}15.1. ci-dessus; alors \var{\inserted[typed]{(la gamme ``géométrique''de 15.2. et la suite}}{(la gamme ``géométrique''de 15.2. et la suite}{insertion}{ega25_25_164-v0014} 15.\unsure{3}. \unsure{à} 15.7. pourraient être séparés
de ces \var{\xout{\unsure{calculs}}}{}{typewriter strike-out}{ega25_25_164-v0015} résultats, en les groupant dans un Nº 17.16, \var{\xout{\illegible{2}}}{}{typewriter strike-out}{ega25_25_164-v0016}
\var{\xout{\unsure{lorsque}}}{}{typewriter strike-out}{ega25_25_164-v0017} De toutes façons, je m'aperçois en cours de rédaction que le
\var{\margin{(}}{\placedmargin{(}}{marginal note placed}{ega25_25_164-v0018}nouveau nº 15 prévu sur les formes lisses éclate en \var{\inserted{au moins}}{au moins}{insertion}{ega25_25_164-v0019} deux numéros assez
distincts et indépendants, l'un contenant des sorites généraux pour
les trucs ``élémentaires de mult n'' , n quelconque, l'autre \var{\deleted{qui}}{}{deletion}{ega25_25_164-v0020} cont\var{\over{ient}{enant}}{enant}{writing over}{ega25_25_164-v0021} la
caractérisation des formes \uline{quadratiques} lisses, et qui n'emprunte au
précédent que le nº 15.1. , i.e. pratiquement la définition de la lissité
d'une forme.

\mspage{164}
